% Generated draft: requires mathematical and source review.
\documentclass[11pt]{article}
\usepackage[margin=25mm]{geometry}
\usepackage{fontspec}
\setmainfont{Latin Modern Roman}
\usepackage{amsmath,amssymb,mathtools,unicode-math}
\setmathfont{Latin Modern Math}
\usepackage{enumitem,needspace,xurl,hyperref}
\setlength{\parindent}{0pt}
\setlength{\parskip}{5pt}
\setlength{\emergencystretch}{3em}
\newcommand{\N}{\mathbb{N}}
\newcommand{\Z}{\mathbb{Z}}
\newcommand{\Q}{\mathbb{Q}}
\newcommand{\R}{\mathbb{R}}
\newcommand{\C}{\mathbb{C}}
\newcommand{\F}{\mathbb{F}}
\newcommand{\A}{\mathbb{A}}
\newcommand{\PP}{\mathbb{P}}
\newcommand{\E}{\mathbb{E}}
\newcommand{\eps}{\varepsilon}
\newcommand{\dd}{\,\mathrm{d}}
\begin{document}
\section*{UnsolvedMath ID: 700772}
\textbf{Source problem code:} catalogue entry
\subsection{Definitions and mathematical statement}
Let \(\Omega\) be a finite set and let \(G\leq\operatorname{Sym}(\Omega)\) be primitive: \(G\) is transitive and preserves no partition of \(\Omega\) other than the partition into singletons and the one-part partition. For \(\alpha\in\Omega\), write
\[
G_\alpha=\{g\in G:g\alpha=\alpha\}.
\]
A suborbit at \(\alpha\) is an orbit of \(G_\alpha\) on \(\Omega\). Suppose that some suborbit \(\Delta\subseteq\Omega\setminus\{\alpha\}\) affords a \(2\)-transitive action of \(G_\alpha\): for any \(x_1,x_2,y_1,y_2\in\Delta\) with \(x_1\ne x_2\) and \(y_1\ne y_2\), there exists \(g\in G_\alpha\) such that \(gx_i=y_i\) for \(i=1,2\). Under the usual nontrivial-degree convention, this also requires \(|\Delta|\geq2\). The supplied source does not explicitly specify that convention; whether singleton suborbits are excluded requires clarification.

A minimal normal subgroup of \(G\) is a nonidentity normal subgroup containing no smaller nonidentity normal subgroup of \(G\). A subgroup \(N\leq G\) is regular on \(\Omega\) if for every \(x,y\in\Omega\) there is exactly one \(n\in N\) with \(nx=y\).

The source attributes to C. E. Praeger (1988) the assertion that every group satisfying the stated hypothesis falls into one of the following cases:
\begin{enumerate}
\item[(a)] \(T\leq G\leq\operatorname{Aut}(T)\) for a finite nonabelian simple group \(T\), with \(T\) identified with its group of inner automorphisms;
\item[(b)] \(G\) has a unique minimal normal subgroup, and that subgroup is regular on \(\Omega\).
\end{enumerate}
\textbf{Requested classification.} In case (a), determine for which classes of finite nonabelian simple groups \(T\) a classification of the primitive permutation groups satisfying the suborbit hypothesis is feasible, and describe the examples as explicitly as possible. In case (b), classify all primitive permutation groups satisfying the suborbit hypothesis whose unique minimal normal subgroup is regular. The source supplies neither a formal criterion for ``feasible'' nor a specified equivalence relation for the classification of permutation actions.

\textbf{Supplied historical qualifications.} The remark dated 1999 reports that two papers of X. G. Fang and C. E. Praeger establish feasibility for Suzuki and Ree simple groups, whereas a paper of J. Wang and C. E. Praeger suggests that feasibility may fail for alternating groups. This last assertion is a reported suggestion, not an impossibility theorem. The remark dated 2001 reports new examples with \(G=\operatorname{PSU}(3,q)\) constructed by X. G. Fang, G. Havas and J. Wang; the record does not specify the parameter restrictions or the permutation actions. The remark dated 2009 reports feasibility in case (a) for sporadic simple groups when the qualifying suborbit is self-paired.

For this last qualification, choose \(\beta\in\Delta\) and let
\[
R=\{(g\alpha,g\beta):g\in G\}.
\]
The suborbit is self-paired when \(R=R^{-1}\), where \(R^{-1}=\{(y,x):(x,y)\in R\}\). In that case \(R\) defines adjacency in an undirected graph on \(\Omega\). The source states that these suborbits correspond to \(2\)-arc-transitive graph actions: \(G\) is transitive on ordered triples \((v_0,v_1,v_2)\) with consecutive vertices adjacent and \(v_0\ne v_2\). This historical result is restricted to self-paired suborbits and does not state a classification for all sporadic cases.
\subsection{Short English statement}
Classify finite primitive permutation groups having a point stabilizer that acts \(2\)-transitively on a suborbit away from the fixed point. The source separates the problem into almost simple groups, where it asks which simple-group families admit an explicit classification, and groups with a unique regular minimal normal subgroup, where it asks for a complete classification. Preserve the reported partial results and their restrictions; ``feasible'' and the treatment of singleton suborbits remain unspecified.
\subsection{Sources}
\begin{itemize}
\item Local UnsolvedMath registry: permanent ID \(700772\), exact source problem code \texttt{catalogue entry}, title \emph{Kourovka Notebook Problem 11.80}. The supplied record labels the problem \texttt{open}; this is a registry status, not an independently verified assessment. Supplied upstream revision: \texttt{d371ee917ac9647fbb06cd6fc529deec6f3dda82}.
\item \emph{The Kourovka Notebook (Unsolved Problems in Group Theory)}, Problem 11.80, proposed by C. E. Praeger, 11th issue (1990); supplied source identifies the 21st issue edition (2026), \texttt{arXiv:1401.0300v48}. The registry describes the collection as edited by E. I. Khukhro and V. D. Mazurov.
\item C. E. Praeger, \emph{J. Austral. Math. Soc. (A)} \textbf{45} (1988), 66--77. Cited by the source for the structural alternatives (a) and (b); no article title is supplied.
\item X. G. Fang and C. E. Praeger, two papers in \emph{Commun. Algebra} \textbf{27} (1999), 3727--3754 and 3755--3769. The source's remark dated 1999 reports feasibility for Suzuki and Ree simple groups.
\item J. Wang and C. E. Praeger, \emph{J. Algebra} \textbf{180} (1996), 808--833. The source's remark dated 1999 describes this paper as suggesting that classification may not be feasible for alternating groups.
\item X. G. Fang, G. Havas and J. Wang, \emph{European J. Combin.} \textbf{20}, no. 6 (1999), 551--557. The source's remark dated 2001 reports new examples with \(G=\operatorname{PSU}(3,q)\).
\item D. Leemans, \emph{J. Algebra} \textbf{322}, no. 3 (2009), 882--892. The source's remark dated 2009 reports feasibility for sporadic simple groups with self-paired \(2\)-transitive suborbits and relates these to \(2\)-arc-transitive actions on undirected graphs.
\item All bibliographic attributions and the remarks dated 1999, 2001 and 2009 are supplied source information, not independently verified literature assessments. No external source was visited for this draft. No explicit reference URLs or separate literature-check date are supplied.
\end{itemize}
\end{document}
